\chapter{持久化、并发与可靠性}

\section{权威数据与身份层级}

当前目标以 PostgreSQL 保存结构化认知权威数据，大型或二进制材料存放于不可变对象库，检索索引按需重建。权威数据、运行时状态和检索视图分别拥有独立身份与生命周期。\cite{nous-research-package}

\begin{ReportWideTable}{p{3.0cm}p{10.2cm}}
\captionof{table}{权威数据、运行时与服务视图的身份层级}\label{tab:identity-layers}
\toprule
\textbf{标识} & \textbf{语义}\\
\midrule
SubjectId & 跨限界上下文保持稳定的逻辑主体身份。\\
object\_id & 稳定的语义对象身份。\\
revision\_id & 不可变内容或解释修订的身份。\\
object\_epoch & 当前版本指针和生命周期状态的并发版本。\\
authority\_seq & 主体范围内有序的提交位置。\\
runtime\_revision & 会话与运行时状态的当前版本。\\
view\_set\_id / generation\_id & 多视图集合与不可变服务代次。\\
request\_id / event\_id & 变更请求与使用事件的幂等标识。\\
\bottomrule
\end{ReportWideTable}

内容修订与抑制、清除等对象状态变化分别处理。时间语义区分事件发生、主体观察、主张有效、系统记录和提交顺序；未知时间与开放区间均需明确表示。

\section{统一变更顺序}

权威数据写入按固定事务顺序完成：

\begin{enumerate}
\item 验证请求、引用、操作者和输入结构；
\item 开启权威数据事务并分配主体提交序号；
\item 按稳定顺序锁定目标对象；
\item 检查幂等键及预期对象版本和当前修订；
\item 重新验证来源、证据、访问控制和生命周期；
\item 写入不可变修订及关联，并更新当前指针、生命周期和版本号；
\item 记录权威数据变更、待发布事件和幂等结果；
\item 提交事务。
\end{enumerate}

模型或网络调用在事务外完成。跨领域更新先由目标领域所有者核验提案，再提交权威数据。\cite{nous-research-package}

\begin{lstlisting}[caption={目标变更合同的简化顺序},label={lst:mutation-order}]
validate(request, actor, source_refs)
BEGIN
  allocate Subject authority_seq
  lock target objects in stable order
  check request idempotency
  check expected epoch and current head
  revalidate evidence, source, ACL and lifecycle
  write immutable revision and relations
  advance head, lifecycle and epoch
  append Authority change and idempotent result
COMMIT
\end{lstlisting}

伪代码概括研究资料包提出的事务语义；各领域实现情况见第十八章。

\section{投影、过期结果与恢复}

权威数据事务按投影类别记录失效事件和处理水位。投影构建器可分别重建精确、词项、稠密、拓扑或摘要视图。发布前核对输入水位，使当前代次始终对应最新有效修订。查询返回前重新核验权威数据、访问控制和生命周期。\cite{nous-research-package,nous-code-snapshot}

请求和事件使用稳定标识处理超时重试、重复投递和乱序交付。重复使用回执按幂等规则去重；同一请求标识对应不同内容摘要时，返回标识冲突。

\section{清除管理}

清除由具有管理权限的流程执行。流程隔离普通读取，取消或标记依赖提案，按规则清理权威数据载荷、服务视图、缓存、上下文和对象库内容，并记录清除回执。需保留的审计元数据只保留不可逆的操作信息，不包含已清除内容、可逆摘要或向量。\cite{nous-research-package}

当前 \SourcePath{crates/memory-service/src/memory_lifecycle.rs} 包含记忆抑制、恢复和清除协调路径。崩溃重试、外部存储与跨领域清除的资格范围需结合相应运行记录确认。\cite{nous-code-snapshot}
